% CHƯƠNG 1. GIỚI THIỆU
\clearpage
\section[Giới thiệu]{GIỚI THIỆU}

Chuỗi thời gian xuất hiện ở khắp nơi: mực nước của một hồ lớn, lưu lượng của một dòng sông, số hành khách hàng không hay chỉ số chứng khoán.

\subsection{Bối cảnh}
Chuỗi thời gian là một dãy quan sát ghi theo thứ tự thời gian. Khác với dữ liệu thông thường, các quan sát liền nhau thường không độc lập: mực nước năm nay phụ thuộc mực nước năm trước, giá hôm nay liên quan đến giá hôm qua. Sự phụ thuộc này vừa là điều cần mô tả, vừa khiến các phương pháp dành cho dữ liệu độc lập không dùng được nguyên vẹn.

ARMA và ARIMA là công cụ nền tảng để mô tả sự phụ thuộc ấy và để dự báo. Lý thuyết của chúng dựa trên những giả thiết gọn, như chuỗi dừng, nhiễu có tính chất cho trước và mô hình đúng bậc, trong khi dữ liệu thực hiếm khi thỏa mãn hoàn hảo vì có thể trôi, đổi mức, có mùa vụ hoặc có biến động thay đổi theo thời gian. Câu hỏi tự nhiên là lý thuyết còn đúng đến đâu trên dữ liệu thực, và khi lệch thì lệch vì lý do gì. Báo cáo này đi tìm lời đáp cho câu hỏi đó.

\subsection{Mục tiêu}
Mục tiêu chung là kiểm chứng bằng Python các kết quả lý thuyết về chuỗi dừng, ARMA, chuỗi không dừng, ARIMA và dự báo trên sáu bộ dữ liệu thực, kết hợp mô phỏng khi cần, rồi đánh giá mức độ phù hợp giữa lý thuyết và thực nghiệm. Cụ thể, báo cáo trả lời ba câu hỏi sau:
\begin{enumerate}[label=\textbf{CH\arabic*:}]
	\item Các kết quả lý thuyết về cấu trúc phụ thuộc và mô hình ARMA có được thỏa mãn trên dữ liệu thực hay không?
	\item Các kết quả lý thuyết về chuỗi không dừng, mô hình ARIMA và dự báo có được thỏa mãn trên dữ liệu thực hay không?
	\item Mức độ phù hợp phụ thuộc thế nào vào đặc điểm của từng bộ dữ liệu, và giả thiết nào bị vi phạm?
\end{enumerate}

\subsection{Phạm vi}
Báo cáo dùng sáu bộ dữ liệu của R, được ký hiệu theo Danh mục ký hiệu và chữ viết tắt ở đầu báo cáo, và chỉ gọi bằng các ký hiệu ấy ở những phần sau. Hai ký hiệu DR và DP cùng chỉ số DAX nhưng ở hai dạng khác nhau, vì suất sinh lợi là dạng dừng còn log giá là dạng không dừng của cùng một chuỗi.
Về phương pháp, báo cáo dùng mô hình ARMA, ARIMA và ARIMA mùa vụ cùng các công cụ đi kèm là tự tương quan, kiểm định Ljung--Box, kiểm định nghiệm đơn vị và kiểm định Diebold--Mariano. Các mô hình như GARCH, Prophet hay học máy nằm ngoài phạm vi, dù có những chuỗi, điển hình là DR, mang đặc điểm mà các mô hình đó xử lý tốt hơn.

Nội dung kiểm chứng gồm 28 kết quả lý thuyết, được liệt kê cùng ký hiệu bộ dữ liệu áp dụng trong Bảng~\ref{tab:danhmuc}. Trong số đó, 20 kết quả được kiểm chứng bằng dữ liệu thực, 7 bằng cả dữ liệu thực lẫn mô phỏng và 1 chỉ bằng mô phỏng.

{\footnotesize
\begin{longtable}{@{}p{2.2cm}p{6.6cm}p{5.2cm}p{2.2cm}@{}}
\caption{Danh mục các nội dung kiểm chứng}\label{tab:danhmuc}\\
\toprule
Địa chỉ & Nội dung & Ký hiệu bộ dữ liệu áp dụng & Hình thức \\
\midrule
\endfirsthead
\toprule
Địa chỉ & Nội dung & Ký hiệu bộ dữ liệu áp dụng & Hình thức \\
\midrule
\endhead
\bottomrule
\endlastfoot
4.1(a) & Tính dừng và nhiễu trắng & HU, LY, DR & Dữ liệu thực \\
4.1(b) & Tính chất của tự hiệp phương sai & HU, LY, DR & Dữ liệu thực \\
4.1(c) & Trung bình mẫu của chuỗi phụ thuộc & HU & Cả hai \\
4.1(d) & Quá trình tuyến tính & HU, LY & Dữ liệu thực \\
4.2(a) & Mô hình AR(1) & HU & Dữ liệu thực \\
4.2(b) & Mô hình AR(p) và hệ Yule–Walker & LY, HU & Dữ liệu thực \\
4.2(c) & Mô hình AR(2) & LY & Dữ liệu thực \\
4.2(d) & Mô hình MA(q) & DR, HU & Cả hai \\
4.2(e) & Mô hình ARMA(p, q) & HU, LY & Dữ liệu thực \\
4.3(a) & Tự tương quan và kiểm định Ljung–Box & HU, LY, DR & Cả hai \\
4.3(b) & Tự tương quan riêng phần & HU, LY, DR & Dữ liệu thực \\
4.3(c) & Ước lượng tham số & LY, HU & Cả hai \\
4.3(d) & Lựa chọn bậc & LY, HU & Dữ liệu thực \\
4.3(e) & Chẩn đoán phần dư & HU, LY, DR & Dữ liệu thực \\
5.1(a) & Xu thế tất định và ngẫu nhiên & NI, AP, DP & Dữ liệu thực \\
5.1(b) & Bước ngẫu nhiên & DP & Mô phỏng \\
5.2(a) & Mô hình ARIMA và độ trôi & NI, DP & Dữ liệu thực \\
5.2(b) & Mô hình ARIMA mùa vụ & AP & Dữ liệu thực \\
5.3(a) & Sai phân và sai phân quá mức & NI, DP & Cả hai \\
5.3(b) & Kiểm định nghiệm đơn vị & NI, AP, DP & Cả hai \\
6.1(a) & Ảnh hưởng vĩnh viễn của cú sốc & NI, DP & Dữ liệu thực \\
6.1(b) & Làm trơn mũ và mô hình mức cục bộ & NI & Dữ liệu thực \\
6.2(a) & Dự báo tối ưu và khoảng dự báo & NI, AP, DP & Cả hai \\
6.2(b) & Dự báo trên thang logarit & AP, DP & Dữ liệu thực \\
6.3(a) & Mốc chuẩn và thước đo sai số & NI, AP, DP & Dữ liệu thực \\
6.3(b) & Gốc cuốn chiếu & NI, AP, DP & Dữ liệu thực \\
6.3(c) & Hàm mất mát và kiểm định Diebold–Mariano & AP, NI & Dữ liệu thực \\
6.3(d) & Quy trình báo cáo dự báo & NI, AP, DP & Dữ liệu thực \\
\end{longtable}}

\subsection{Bố cục}
Báo cáo gồm tám chương. Chương 2 trình bày cơ sở lý thuyết theo ba mục là chuỗi dừng và ARMA, chuỗi không dừng và ARIMA, dự báo và đánh giá dự báo. Chương 3 giới thiệu dữ liệu và phương pháp. Ba mục của Chương 2 lần lượt được kiểm chứng ở các Chương 4, 5 và 6, sau đó Chương 7 thảo luận chung và Chương 8 kết luận.

Mỗi kết quả ở Chương 2 mang một địa chỉ dạng \texttt{2.m.n(z)} và được kiểm chứng tại mục có địa chỉ \texttt{(m+3).n(z)}; chẳng hạn kết quả \texttt{2.1.2(c)} được kiểm chứng tại \texttt{4.2(c)}. Mỗi nội dung kiểm chứng gồm sáu phần: kết quả cần kiểm chứng, cách triển khai, mã Python, kết quả tính toán, phân tích và diễn giải, và nhận xét. Phần nhận xét xếp loại kết quả thành \emph{phù hợp}, \emph{phù hợp một phần} hoặc \emph{không phù hợp} theo thang trong Bảng~\ref{tab:thangdo}.
